% Body of the Time-series report — used by timeseries.tex (as an article) and G3-report.tex (as a chapter).
% Every number comes from generated/timeseries.tex (built from timeseries/results.json).

\section{Problem and objective}\label{ts:problem}
Hanoi has a real winter: cold surges of the north-east monsoon can cut the temperature by 5–10~°C within a day, and those
are exactly the days when a forecast matters most for health, energy use and agriculture.

\textbf{Objective.} Forecast \textbf{tomorrow's daily mean temperature in Hanoi} from today's and past weather, then measure
how quickly forecast skill decays when predicting up to seven days ahead.

\textbf{How success is measured.} Mean absolute error (MAE, °C) on a chronologically held-out test period, and \emph{skill}
relative to simple baselines — above all \emph{persistence} (``tomorrow = today''), which is hard to beat for one-day forecasts.

\section{Dataset}\label{ts:dataset}
The data was downloaded by the group from the \textbf{Open-Meteo Historical Weather API} \cite{ts-openmeteo} (free, no API
key, CC BY 4.0), which serves the \textbf{ERA5 reanalysis} \cite{ts-era5}: a physically consistent reconstruction of past
weather on a 9–25~km grid, combining a weather model with millions of observations. The exact request is included in the
notebook (\code{fetch\_open\_meteo}); \emph{Run all} uses the saved CSV so results are reproducible.

\begin{itemize}
\item Location: central Hanoi (21.03~°N, 105.85~°E; nearest grid cell 21.05~°N, 105.90~°E, elevation 19~m).
\item Period: \TSStart{} – \TSEnd{}, daily, local time (GMT+7): \TSDays{} consecutive days, \TSVars{} variables.
\end{itemize}

\begin{table}[H]\centering\small
\caption{Variables. The target is the daily mean temperature.}\label{ts:tab-vars}
\begin{tabularx}{\linewidth}{L{6.4cm}lY}
\toprule \hd{Variable} & \hd{Unit} & \hd{Meaning} \\ \midrule
\code{temperature\_2m\_mean} / \code{\_max} / \code{\_min} & °C & daily mean (\textbf{target}), maximum, minimum at 2~m \\
\code{apparent\_temperature\_mean} & °C & ``feels-like'' temperature \\
\code{precipitation\_sum}, \code{precipitation\_hours} & mm, h & daily rain and rainy hours \\
\code{relative\_humidity\_2m\_mean}, \code{dew\_point\_2m\_mean} & \%, °C & moisture \\
\code{pressure\_msl\_mean} & hPa & mean sea-level pressure \\
\code{cloud\_cover\_mean} & \% & cloud cover \\
\code{wind\_speed\_10m\_max}, \code{wind\_direction\_10m\_dominant} & km/h, ° & maximum wind speed, dominant direction \\
\code{shortwave\_radiation\_sum}, \code{sunshine\_duration} & MJ/m², s & solar energy, sunshine \\
\bottomrule
\end{tabularx}
\end{table}

\begin{table}[H]\centering\small
\caption{General requirements (section 4.1 and 5) and how this dataset meets them.}\label{ts:tab-req}
\begin{tabularx}{\linewidth}{L{4.2cm}Yl}
\toprule \hd{Requirement} & \hd{This dataset} & \hd{Status} \\ \midrule
Enough samples & \TSDays{} days & \met \\
Numerical columns (scaling) & \TSVars{} weather variables plus engineered lags & \met \\
Categorical columns (encoding) & derived: season and 8-sector wind direction (one-hot), circular encodings & \met \\
Missing values handled & none in the source (verified); gaps simulated to compare imputation methods (Section~\ref{ts:missing}) & \met \\
Time-aware split & train / validation / test by date; time-series cross-validation & \met \\
\bottomrule
\end{tabularx}
\end{table}

\section{Exploratory data analysis}\label{ts:eda}
\fig{ts_02_seasonality}{Left: climatology by day of year (mean, ±1 standard deviation, record range). Right: mean absolute day-to-day change by month.\label{ts:fig-season}}
\begin{finding}
\Q How different are the seasons — and when is the weather hardest to predict?
\F Summer (June–August) averages \textbf{\TSSummer~°C} and winter (December–February) \textbf{\TSWinter~°C}, a swing of about
\TSSeasonGap~°C. Winter is also far more \emph{variable}: the standard deviation within January is \TSJanStd~°C versus
\TSJulStd~°C in July, and the average day-to-day change peaks in February–March (\TSDtdFeb–\TSDtdMar~°C) against \TSDtdJul~°C
in July. Forecasting should therefore be hardest in winter and spring.
\end{finding}

\fig{ts_03_trend}{Left: annual mean temperature with a linear trend. Right: monthly anomaly versus the 2005–2024 average (blue = colder, red = warmer).\label{ts:fig-trend}}
\begin{finding}
\Q Is Hanoi getting warmer?
\F The annual mean rises by \textbf{\TSTrend~°C per decade} over 2005–2025; 2021–25 averaged \TSLastFive~°C against
\TSFirstFive~°C in 2005–09 (\TSDeltaFive~°C). Large cold anomalies cluster in the early years (February 2008, January–February
2011) and recent years are mostly warm. Twenty-one years is short and year-to-year swings are large (\TSAnnualMin~°C in
\TSAnnualMinYear{} vs \TSAnnualMax~°C in \TSAnnualMaxYear), so the slope is uncertain — but it matters for modelling: a
climatology learned on 2005–2021 will be slightly too cold for 2024–2026 (Section~\ref{ts:results}).
\end{finding}

\begin{figure}[H]\centering
\begin{subfigure}[t]{0.58\linewidth}\centering\includegraphics[width=\linewidth]{ts_04_acf}
\caption{Autocorrelation of temperature and of its anomaly.}\end{subfigure}\hfill
\begin{subfigure}[t]{0.4\linewidth}\centering\includegraphics[width=\linewidth]{ts_05_cold_surge}
\caption{Pressure change today vs temperature change tomorrow (Nov–Mar).}\end{subfigure}
\caption{Memory of the series and the cold-surge signal.\label{ts:fig-memory}}
\end{figure}
\begin{finding}
\Q How much does today tell us about tomorrow, and can we see a cold surge coming?
\F The raw series is dominated by the yearly cycle (autocorrelation \TSAcfYear{} at lag 365). Once the season is removed, the
anomaly still has strong day-to-day memory (\textbf{\TSAcfAnomOne{} at lag 1}) that fades to \textbf{\TSAcfAnomSeven{} by lag 7}:
a warm or cold spell ``remembers'' itself for only a few days. In the cold season a pressure rise today warns of cooling
tomorrow (r = \TSR): on the \TSNRise{} days when pressure rose by more than 3~hPa, the next day changed by \TSDTRise~°C on
average (vs \TSDTOther~°C otherwise) and a drop of more than 2~°C was \textbf{\TSPRatio× as likely} (\TSPRise{} vs \TSPOther).
\end{finding}

\section{Data preparation}\label{ts:prep}
\subsection{Continuity and missing values}\label{ts:missing}
A time series must be checked for missing \emph{dates}, not only missing cells, because a skipped day silently shifts every
lag feature. The file has \TSDays{} rows for \TSDays{} calendar days, no duplicated dates and \TSMissingDates{} missing cells.
Real station data is rarely this clean, so we used the complete series as ground truth for an experiment: hide random blocks of
1–30 consecutive days (about 5\% of the series, 5 repetitions) and reconstruct them with four methods.

\begin{table}[H]\centering\small
\caption{Imputation error (MAE in °C on the hidden days; best per gap length in bold).}\label{ts:tab-imputation}
% Generated by tools/build_reports.py — do not edit.
\begin{tabular}{lrrrr}
\toprule
\hd{Gap length} & \hd{Forward fill} & \hd{Linear interp.} & \hd{Climatology} & \hd{Seasonal interp.} \\
\midrule
1 day & 1.16 & 0.71 & 1.97 & \textbf{0.70} \\
3 days & 1.76 & 1.14 & 1.89 & \textbf{1.13} \\
7 days & 2.21 & 1.69 & 2.00 & \textbf{1.63} \\
14 days & 2.45 & 1.95 & \textbf{1.86} & 1.93 \\
30 days & 2.80 & 2.22 & \textbf{1.87} & 2.15 \\
\bottomrule
\end{tabular}

\end{table}
\begin{finding}
\Q If a station lost data for a few days, how should the gap be filled — and does the answer depend on its length?
\F For short gaps, \textbf{seasonal interpolation} (interpolating the anomaly, then adding the climatology back) is best:
\TSImpSeasOne, \TSImpSeasThree{} and \TSImpSeasSeven~°C for 1, 3 and 7 days, slightly ahead of plain linear interpolation
(\TSImpLinOne, \TSImpLinThree, \TSImpLinSeven). Forward fill is worst at every length. For gaps of \textbf{14 days or more,
the seasonal normal alone wins} (\TSImpClimFourteen~°C), because the anomaly's memory (about a week) has run out. Rule:
interpolate the anomaly up to about a week, fall back to climatology for longer outages.
\end{finding}

\subsection{Outliers}
Temperature outliers must be judged against the season (12~°C is normal in January but extreme in July), so each day's anomaly
is divided by the standard deviation for that day of the year and |z| > 3 is flagged. Only \TSExtremes{} days qualify, all
cold (March 2011, early October 2011, September 2013): real cold spells, exactly what a forecast must handle, so they are
\textbf{kept}. The record \TSColdest~°C of \TSColdestDate{} is not even flagged because January is so variable.
Precipitation has \TSRainOutliers{} IQR outliers and skewness \TSRainSkew{} (maximum \TSRainMax~mm on \TSRainMaxDate); heavy
rain is also real, so it is kept and fed to the models as \code{log1p(rain)} (skewness \TSRainSkewLog).

\subsection{Features, encoding and scaling}
For a forecast issued at the end of day \emph{t} for day \emph{t + h}, every feature must be known on day \emph{t}. The class
\code{WeatherFeatureBuilder} builds them; its climatology is learned from the \textbf{training years only}.

\begin{table}[H]\centering\small
\caption{Features (30 in total) and their encoding.}\label{ts:tab-features}
\begin{tabularx}{\linewidth}{L{3.5cm}YL{3.4cm}}
\toprule \hd{Group} & \hd{Features} & \hd{Encoding / scaling} \\ \midrule
Recent temperature & today; lags 1, 2, 3, 7, 14 days; rolling means over 3, 7, 30 days; today's max, min and range & standardised \\
Anomaly & today − climatology (smoothed day-of-year mean) & standardised \\
Weather today & humidity, dew point, pressure, pressure change over 1 and 3 days, cloud cover, wind speed, rain, radiation, sunshine & \code{log1p} for rain; standardised \\
Wind direction & sine and cosine of the direction & circular (359° is next to 1°) \\
Calendar of target day & day of year; climatology of the target day & sine/cosine (cyclical) \\
Categorical & season of target day (DJF, MAM, JJA, SON); 8-sector wind direction (N, NE, …, NW) & one-hot \\
\bottomrule
\end{tabularx}
\end{table}
The first 30 days (rolling-window warm-up) and the last day (no target) are dropped. Scaling uses \code{StandardScaler} inside
the pipeline; it is required by Ridge regression and harmless for tree models.

\subsection{Time-based split}
A random split would let the model learn from days \emph{after} the ones it is tested on. We split by the date being forecast:
\textbf{train} 2005–2021 (\TSTrain{} days), \textbf{validation} 2022–2023 (\TSVal{} days) and \textbf{test} 1~January~2024 –
\TSEnd{} (\TSTest{} days). Hyperparameters are tuned with \code{TimeSeriesSplit} (5 expanding windows inside the training years),
which always trains on the past and validates on the following period.

\section{Modelling}\label{ts:models}
\textbf{Baselines.} \emph{Persistence} (tomorrow = today); \emph{climatology} (the average for that calendar day); and
\emph{damped anomaly persistence} (climatology + φ × today's anomaly, with φ = \TSPhi, the lag-1 autocorrelation of anomalies
in the training years).

\textbf{Machine-learning models.} Ridge regression, Random Forest and HistGradientBoosting (HGB), all in the same
pipeline. HGB is trained either on the temperature itself or on the \textbf{change} ΔT = T(t+1) − T(t), adding today's value
back afterwards (\code{DeltaRegressor} class). Trees cannot extrapolate beyond the range seen in training, and a warming climate
makes that a real risk; predicting the change avoids it.

\textbf{Tuning.} A random search over 20 HGB configurations (scikit-learn defaults included) with time-series cross-validation
preferred a more regularised model (CV MAE \TSCvTuned{} vs \TSCvDefault{} for the defaults), but it did \emph{worse} on
2022–2023 (\TSValTuned{} vs \TSValFinal). A 0.015~°C gain on older folds is within year-to-year noise, so the final model is
the one chosen on validation: HGB on the change, with default hyperparameters.

\begin{table}[H]\centering\small
\caption{One-day-ahead forecasts: validation (2022–2023) and test (2024 – September 2026) scores. Skill = 1 − MAE / MAE of persistence (test).}\label{ts:tab-models}
% Generated by tools/build_reports.py — do not edit.
\begin{tabular}{lrrrrr}
\toprule
\hd{Model} & \hd{Val MAE} & \hd{Test MAE} & \hd{Test RMSE} & \hd{Test R²} & \hd{Skill} \\
\midrule
Climatology (seasonal normal) & 2.014 & 1.990 & 2.538 & 0.722 & −84.2\% \\
Persistence (tomorrow = today) & 1.105 & 1.081 & 1.504 & 0.902 & 0.0\% \\
Damped anomaly persistence & 1.092 & 1.075 & 1.434 & 0.911 & 0.5\% \\
Ridge regression & 0.974 & 0.970 & 1.283 & 0.929 & 10.2\% \\
Random Forest & 0.921 & 0.935 & 1.245 & 0.933 & 13.4\% \\
HGB, temperature as target & 0.933 & 0.932 & 1.231 & 0.934 & 13.8\% \\
HGB, change as target, tuned & 0.922 & 0.919 & 1.222 & 0.935 & 15.0\% \\
\textbf{HGB, change as target (final)} & \textbf{0.906} & \textbf{0.916} & \textbf{1.226} & \textbf{0.935} & \textbf{15.3\%} \\
\bottomrule
\end{tabular}

\end{table}
\begin{finding}
\Q Do machine-learning models beat the classical baselines, and which framing works best?
\F Every model beats every baseline. HGB predicting the change is best on validation (\TSValFinal~°C, \TSValSkill{} below
persistence at \TSValPersist) and beats HGB predicting the level (\TSValLevel). Ridge regression is weakest (\TSValRidge): the
useful signals interact non-linearly — a northerly wind matters mostly in winter. Damped anomaly persistence barely improves on
persistence (\TSValDamped) and climatology alone is poor (\TSValClim~°C): for \emph{tomorrow}, today's weather matters far more
than the calendar.
\end{finding}

\section{Results on the test set}\label{ts:results}
On the \TSTest{} unseen days the final model reaches \textbf{MAE \TSTestFinal~°C, RMSE \TSTestRmse~°C and R² \TSTestRtwo},
a \textbf{\TSTestSkill{} lower error than persistence} (\TSTestPersist~°C); the ranking of all models matches validation
(Table~\ref{ts:tab-models}).

\fig[0.86\linewidth]{ts_10_forecast_vs_actual}{Forecast versus actual over the whole test period (top) and the hardest winter window (bottom), with persistence for comparison.\label{ts:fig-forecast}}

\begin{table}[H]\centering\small
\caption{Test MAE by month (bold = final model better than persistence).}\label{ts:tab-months}
\footnotesize\setlength{\tabcolsep}{3.6pt}
% Generated by tools/build_reports.py — do not edit.
\begin{tabular}{lrrrrrrrrrrrr}
\toprule
\hd{Test MAE (°C)} & \hd{Jan} & \hd{Feb} & \hd{Mar} & \hd{Apr} & \hd{May} & \hd{Jun} & \hd{Jul} & \hd{Aug} & \hd{Sep} & \hd{Oct} & \hd{Nov} & \hd{Dec} \\
\midrule
Persistence & 1.36 & 1.61 & 1.24 & 1.08 & 1.08 & 0.99 & 0.75 & 0.87 & 0.86 & 0.96 & 1.01 & 1.12 \\
Final model & \textbf{1.05} & \textbf{1.16} & \textbf{0.82} & 1.10 & \textbf{0.91} & \textbf{0.87} & \textbf{0.69} & \textbf{0.84} & \textbf{0.72} & \textbf{0.90} & 1.12 & \textbf{0.88} \\
\bottomrule
\end{tabular}

\end{table}

\begin{table}[H]\centering\small
\caption{The \TSWorstN{} largest test errors (°C; pressure change on the forecast day in hPa).}\label{ts:tab-worst}
% Generated by tools/build_reports.py — do not edit.
\begin{tabular}{lrrrrr}
\toprule
\hd{Date} & \hd{Actual} & \hd{Forecast} & \hd{Error} & \hd{1-day change} & \hd{Pressure change} \\
\midrule
29 March 2025 & 17.0 & 23.7 & +6.7 & −8.5 & +4.2 \\
8 February 2026 & 15.0 & 20.7 & +5.7 & −7.0 & +2.2 \\
18 November 2025 & 17.8 & 22.7 & +4.9 & −4.9 & −0.1 \\
23 February 2024 & 19.3 & 23.8 & +4.5 & −6.0 & +3.7 \\
26 January 2025 & 14.9 & 19.4 & +4.5 & −6.4 & +1.6 \\
29 April 2024 & 31.7 & 27.4 & −4.3 & +1.0 & +2.3 \\
7 February 2024 & 15.7 & 20.0 & +4.3 & −5.7 & +0.8 \\
22 January 2024 & 11.3 & 15.5 & +4.2 & −4.6 & +3.5 \\
\bottomrule
\end{tabular}

\end{table}
\begin{finding}
\Q When does the model fail?
\F The biggest gains over persistence come in winter and spring — February \TSMoFebModel{} vs \TSMoFebPersist, March
\TSMoMarModel{} vs \TSMoMarPersist, December \TSMoDecModel{} vs \TSMoDecPersist~°C — while the model is slightly worse in the
transition months (April \TSMoAprModel{} vs \TSMoAprPersist, November \TSMoNovModel{} vs \TSMoNovPersist). \textbf{\TSWorstCold{} of
the \TSWorstN{} worst forecasts are cold surges}: the temperature fell \TSWorstDropMin–\TSWorstDropMax~°C in one day and the
forecast was \TSWorstErrMin–\TSWorstErrMax~°C too warm, usually after a pressure rise that the model reacted to only partly. On
average the residual is \TSResMean~°C (standard deviation \TSResStd): forecasts are slightly too cold, consistent with the
warming of Figure~\ref{ts:fig-trend}.
\end{finding}

\section{Extra experiments}\label{ts:extra}
\subsection{How far ahead is the model useful?}
The features are rebuilt and the final configuration retrained for horizons of 1 to 7 days; all methods are scored on the same
test period.

\begin{table}[H]\centering\small
\caption{Test MAE (°C) by forecast horizon (best per horizon in bold).}\label{ts:tab-horizon}
% Generated by tools/build_reports.py — do not edit.
\begin{tabular}{lrrrrrrr}
\toprule
\hd{Test MAE (°C)} & \hd{1 d} & \hd{2 d} & \hd{3 d} & \hd{4 d} & \hd{5 d} & \hd{6 d} & \hd{7 d} \\
\midrule
Persistence (tomorrow = today) & 1.08 & 1.74 & 2.16 & 2.44 & 2.56 & 2.60 & 2.61 \\
Climatology (seasonal normal) & 1.99 & 1.99 & 1.99 & 1.99 & 1.99 & \textbf{1.99} & 1.99 \\
Damped anomaly persistence & 1.07 & 1.60 & 1.86 & 1.97 & 2.00 & 1.99 & \textbf{1.99} \\
Final model & \textbf{0.92} & \textbf{1.49} & \textbf{1.79} & \textbf{1.95} & \textbf{1.99} & 2.00 & 1.99 \\
\bottomrule
\end{tabular}

\end{table}
\fig[0.78\linewidth]{ts_13_horizon}{Forecast error versus horizon on the test period.\label{ts:fig-horizon}}
\begin{finding}
\Q At what horizon does the model stop beating simple baselines?
\F The model beats every baseline for \textbf{1 to \TSUseful{} days ahead} (MAE \TSHOne, \TSHTwo, \TSHThree, \TSHFour~°C vs
\TSHOneBase, \TSHTwoBase, \TSHThreeBase, \TSHFourBase{} for the best baseline). From day 5 all methods converge on climatology
(about \TSClim~°C): the anomaly memory has faded, exactly as the autocorrelation of Figure~\ref{ts:fig-memory} predicted.
Persistence becomes worse than the seasonal average from day \TSPersistCross{} (\TSPersistCrossMae~°C).
\end{finding}

\subsection{What drives the predicted change?}
\fig[0.78\linewidth]{ts_12_importance}{Permutation importance of each feature for the predicted one-day change (test set).\label{ts:fig-importance}}
\begin{finding}
\Q Which signals does the model use to predict tomorrow's change?
\F Today's \textbf{anomaly} (\TSImpAnom) comes first — an unusually warm day tends to cool back towards normal and vice versa.
Next is the \textbf{north–south component of the wind} (\TSImpWind): a northerly wind brings cold air, the signature of a cold
surge. Rain (\TSImpRain) and cloud cover (\TSImpCloud) follow, as they suppress daytime heating. The one-day pressure change is
used (\TSImpDp) but ranks lower than Figure~\ref{ts:fig-memory} suggests, because the wind direction carries much of the same
information.
\end{finding}

\section{Conclusions}\label{ts:conclusion}
\begin{enumerate}
\item \textbf{Machine learning beats the classical baselines, modestly}: \TSTestSkill{} lower test MAE than persistence
(\TSTestFinal{} vs \TSTestPersist~°C). Framing the target as the \emph{change} helped more than hyperparameter tuning.
\item \textbf{The hard days are cold surges.} Winter and spring are where the model gains most and where its worst misses happen;
the physical signals — northerly wind, rising pressure, today's anomaly — are what it learned to use.
\item \textbf{Skill fades after about \TSUseful{} days}, when every method converges on the seasonal average, matching the
roughly one-week memory of temperature anomalies.
\item \textbf{Data preparation lessons}: split by time, never randomly; fit the climatology on training years only; fill short
gaps by interpolating the anomaly and long gaps with climatology.
\item \textbf{Hanoi is warming} by about \TSTrend~°C per decade in this record — enough to leave a \TSResMean~°C bias in
forecasts trained on earlier years.
\end{enumerate}
\textbf{Limitations.} ERA5 is a smoothed grid-cell value rather than a thermometer in the city; real station data would be
noisier. The forecast uses the complete weather of day \emph{t}, slightly more than an evening forecast would have. Only local
variables are used, while cold surges start hundreds of kilometres to the north.

\textbf{Future work.} Add upstream features (pressure and temperature in southern China a day earlier) to see cold surges
coming; use hourly data for 24-hour forecasts; produce probabilistic forecasts (quantile regression) to warn of large drops;
compare with archived forecasts from a physics-based weather model; retrain regularly to follow the warming trend.

\begin{thebibliography}{9}
\bibitem{ts-openmeteo} Open-Meteo. \emph{Historical Weather API} (weather data by Open-Meteo.com, CC BY 4.0).
  \url{https://open-meteo.com/en/docs/historical-weather-api}
\bibitem{ts-era5} H.~Hersbach et al. \emph{The ERA5 global reanalysis}. Quarterly Journal of the Royal Meteorological Society,
  146 (730), 2020, 1999--2049. \url{https://doi.org/10.1002/qj.3803}
\bibitem{ts-sklearn} F.~Pedregosa et al. \emph{Scikit-learn: Machine Learning in Python}. Journal of Machine Learning
  Research, 12 (2011), 2825--2830.
\bibitem{ts-nb} Group G3. \emph{Notebook} \code{G3\_TimeSeries\_HanoiWeather.ipynb}. \TSNotebook
\bibitem{ts-page} Group G3. \emph{Project page (rendered notebook, report and video)}. \TSPage\ — video: \TSVideo
\end{thebibliography}
